\documentclass[aspectratio=169,11pt]{beamer}
\usetheme{Madrid}
\usecolortheme{default}
\usepackage[utf8]{inputenc}
\usepackage[T1]{fontenc}
\usepackage[spanish]{babel}
\usepackage{lmodern}
\usepackage{booktabs}
\usepackage{array}
\usepackage{xcolor}
\usepackage{tikz}
\usetikzlibrary{arrows.meta,positioning,shapes.rounded}

\definecolor{azul}{RGB}{8,120,201}
\definecolor{verde}{RGB}{88,204,2}
\definecolor{gris}{RGB}{96,112,128}
\definecolor{fondo}{RGB}{246,251,255}
\setbeamercolor{structure}{fg=azul}
\setbeamercolor{frametitle}{fg=white,bg=azul}
\setbeamercolor{title}{fg=white,bg=azul}
\setbeamertemplate{navigation symbols}{}
\setbeamertemplate{footline}[frame number]

\title[Sistema HTML Seguro]{Sistema HTML Seguro con Google Apps Script}
\subtitle{Carga, activación, usuarios, roles e informes con una interfaz de clics}
\author{Tutorial de configuración y uso}
\date{Versión 1.0}

\newcommand{\boton}[1]{\textcolor{azul}{\textbf{#1}}}
\newcommand{\ruta}[1]{\texttt{#1}}
\newcommand{\alertbox}[1]{\begin{block}{Importante}#1\end{block}}

\begin{document}

\begin{frame}
  \titlepage
\end{frame}

\begin{frame}{Objetivo del sistema}
\begin{itemize}
  \item Administrar prácticas HTML sin exigir conocimientos técnicos al personal docente.
  \item Activar o desactivar aplicaciones mediante interruptores.
  \item Controlar acceso mediante usuario, contraseña y rol.
  \item Permitir que estudiantes entren solo a las prácticas autorizadas.
  \item Registrar accesos e informes de uso.
\end{itemize}
\alertbox{Después de la instalación inicial, la administración cotidiana se realiza desde una página web con botones.}
\end{frame}

\begin{frame}{Arquitectura general}
\centering
\begin{tikzpicture}[node distance=1.3cm, every node/.style={font=\small}]
\tikzstyle{box}=[rounded corners, draw=azul, thick, fill=fondo, align=center, minimum width=3.0cm, minimum height=.8cm]
\node[box] (admin) {Administrador / Docente\\Panel de control};
\node[box, right=of admin] (gas) {Google Apps Script\\Web App};
\node[box, right=of gas] (sheet) {Google Sheet\\usuarios, estados, informes};
\node[box, below=of sheet] (drive) {Google Drive\\HTML privados};
\node[box, below=of admin] (student) {Estudiante\\Práctica autorizada};
\draw[-{Latex}, thick] (admin) -- (gas);
\draw[-{Latex}, thick] (student) -- (gas);
\draw[-{Latex}, thick] (gas) -- (sheet);
\draw[-{Latex}, thick] (gas) -- (drive);
\end{tikzpicture}
\vspace{.4cm}
\small El usuario no toca JSON, GitHub, tokens ni configuraciones avanzadas. Todo se guarda en Google Sheets y Drive.
\end{frame}

\begin{frame}{Roles disponibles}
\begin{tabular}{p{3cm}p{9.5cm}}
\toprule
\textbf{Rol} & \textbf{Permisos principales}\\
\midrule
ADMINISTRADOR & Crea usuarios, carga HTML, elimina aplicaciones, activa/desactiva, programa horarios, ve informes.\\
DOCENTE & Activa/desactiva aplicaciones, programa horarios y ve informes. No administra usuarios ni carga HTML.\\
ESTUDIANTE & Ingresa a prácticas autorizadas y genera/consulta sus propios informes.\\
\bottomrule
\end{tabular}
\end{frame}

\begin{frame}{Archivos incluidos en el comprimido}
\begin{block}{Carpeta principal}
\begin{itemize}
  \item \ruta{apps\_script/Code.gs}: lógica del servidor.
  \item \ruta{apps\_script/Index.html}: interfaz completa.
  \item \ruta{apps\_script/appsscript.json}: configuración del proyecto.
  \item \ruta{html\_para\_subir/}: prácticas HTML de ejemplo.
  \item \ruta{documentacion/Tutorial\_Sistema\_HTML\_Seguro.pdf}: este tutorial.
\end{itemize}
\end{block}
\alertbox{No edite el código salvo que quiera personalizar el sistema. La operación cotidiana se hace desde el panel.}
\end{frame}

\begin{frame}{Requisitos previos}
\begin{itemize}
  \item Una cuenta Google.
  \item Acceso a Google Sheets y Google Apps Script.
  \item Los archivos del comprimido descomprimidos en su computador.
  \item Los HTML que desea cargar como prácticas.
\end{itemize}
\begin{block}{No se requiere}
\begin{itemize}
  \item Servidor propio.
  \item GitHub.
  \item Cloudflare.
  \item Tarjeta de crédito.
  \item Conocimientos de administración web para uso diario.
\end{itemize}
\end{block}
\end{frame}

\begin{frame}{Paso 1: crear la hoja base}
\begin{enumerate}
  \item Abra Google Sheets.
  \item Cree una hoja nueva.
  \item Asigne un nombre claro, por ejemplo:\\
  \ruta{Sistema HTML Seguro - Control de prácticas}
  \item No cree pestañas manualmente.
\end{enumerate}
\alertbox{El propio script creará las pestañas necesarias: USUARIOS, APLICACIONES, SESIONES, INFORMES y AUDITORIA.}
\end{frame}

\begin{frame}{Paso 2: abrir Apps Script}
\begin{enumerate}
  \item Dentro de la hoja, vaya a \boton{Extensiones}.
  \item Seleccione \boton{Apps Script}.
  \item Se abrirá un proyecto asociado a esa hoja.
  \item Cambie el nombre del proyecto si lo desea.
\end{enumerate}
\begin{block}{Recomendación}
Use un nombre fácil de reconocer, por ejemplo \ruta{SistemaHTMLSeguro}.
\end{block}
\end{frame}

\begin{frame}{Paso 3: copiar Code.gs}
\begin{enumerate}
  \item En Apps Script, abra el archivo \ruta{Code.gs} existente.
  \item Borre su contenido.
  \item Abra el archivo del paquete:\\
  \ruta{apps\_script/Code.gs}
  \item Copie todo su contenido y péguelo en Apps Script.
  \item Guarde el proyecto.
\end{enumerate}
\end{frame}

\begin{frame}{Paso 4: crear Index.html}
\begin{enumerate}
  \item En Apps Script, haga clic en \boton{+} junto a \boton{Archivos}.
  \item Seleccione \boton{HTML}.
  \item Nombre el archivo exactamente:\\
  \ruta{Index}
  \item Abra \ruta{apps\_script/Index.html} del paquete.
  \item Copie todo el contenido y péguelo en el archivo \ruta{Index.html}.
  \item Guarde.
\end{enumerate}
\alertbox{El nombre debe ser \ruta{Index}; Apps Script añadirá la extensión .html automáticamente.}
\end{frame}

\begin{frame}{Paso 5: pegar appsscript.json}
\begin{enumerate}
  \item En Apps Script, abra \boton{Configuración del proyecto}.
  \item Active \boton{Mostrar archivo de manifiesto appsscript.json}.
  \item Vuelva al editor.
  \item Abra \ruta{appsscript.json}.
  \item Reemplace su contenido con el archivo:\\
  \ruta{apps\_script/appsscript.json}
  \item Guarde.
\end{enumerate}
\end{frame}

\begin{frame}{Paso 6: autorizar permisos}
\begin{enumerate}
  \item En Apps Script, seleccione la función \ruta{doGet}.
  \item Pulse \boton{Ejecutar}.
  \item Google pedirá autorización.
  \item Revise los permisos y acepte con la cuenta propietaria.
\end{enumerate}
\begin{block}{Permisos usados}
\begin{itemize}
  \item Google Sheets: guardar usuarios, estados e informes.
  \item Google Drive: guardar los archivos HTML privados.
\end{itemize}
\end{block}
\end{frame}

\begin{frame}{Paso 7: desplegar como aplicación web}
\begin{enumerate}
  \item Pulse \boton{Implementar} o \boton{Deploy}.
  \item Seleccione \boton{Nueva implementación}.
  \item En tipo, elija \boton{Aplicación web}.
  \item Configure:
  \begin{itemize}
    \item Ejecutar como: \boton{Yo}.
    \item Quién tiene acceso: \boton{Cualquier persona con el enlace}.
  \end{itemize}
  \item Pulse \boton{Implementar}.
  \item Copie la URL de la aplicación web.
\end{enumerate}
\alertbox{Esa URL será el panel único para administradores, docentes y estudiantes.}
\end{frame}

\begin{frame}{Paso 8: crear el primer ADMINISTRADOR}
\begin{enumerate}
  \item Abra la URL de la aplicación web.
  \item Si no existen usuarios, aparecerá la pantalla de primer uso.
  \item Cree el usuario ADMINISTRADOR inicial.
  \item Use una contraseña de al menos 10 caracteres.
\end{enumerate}
\begin{block}{Sugerencia}
Anote esta cuenta en un lugar seguro. Será la cuenta que puede crear docentes y estudiantes.
\end{block}
\end{frame}

\begin{frame}{Paso 9: cargar aplicaciones HTML}
\begin{enumerate}
  \item Ingrese como ADMINISTRADOR.
  \item Abra la pestaña \boton{Cargar HTML}.
  \item Escriba un título visible.
  \item Seleccione un archivo .html.
  \item Pulse \boton{Subir aplicación}.
\end{enumerate}
\alertbox{Toda aplicación nueva queda DESACTIVADA por defecto. Esto evita publicarla accidentalmente.}
\end{frame}

\begin{frame}{Paso 10: activar o desactivar una práctica}
\begin{enumerate}
  \item Vaya a la pestaña \boton{Aplicaciones}.
  \item Busque la práctica.
  \item Cambie el interruptor \boton{Activada}.
  \item El cambio queda guardado en la hoja central.
\end{enumerate}
\begin{block}{Quién puede hacerlo}
\begin{itemize}
  \item ADMINISTRADOR: sí.
  \item DOCENTE: sí.
  \item ESTUDIANTE: no.
\end{itemize}
\end{block}
\end{frame}

\begin{frame}{Paso 11: programar apertura y cierre}
\begin{enumerate}
  \item En cada aplicación verá campos \boton{Inicio} y \boton{Fin}.
  \item Defina fecha y hora de apertura.
  \item Defina fecha y hora de cierre.
  \item La disponibilidad se calcula automáticamente.
\end{enumerate}
\begin{block}{Ejemplo}
Una práctica puede estar activa, pero disponible solo entre 08:00 y 10:00.
\end{block}
\end{frame}

\begin{frame}{Paso 12: crear usuarios}
\begin{enumerate}
  \item Ingrese como ADMINISTRADOR.
  \item Abra la pestaña \boton{Usuarios}.
  \item Pulse \boton{Nuevo usuario}.
  \item Complete usuario, nombre, rol y contraseña.
  \item Seleccione prácticas permitidas.
  \item Guarde.
\end{enumerate}
\alertbox{Para ESTUDIANTE, marque solo las prácticas a las que debe entrar. Puede usar Todas (*) si corresponde.}
\end{frame}

\begin{frame}{Uso del DOCENTE}
\begin{itemize}
  \item Ingresa con su usuario y contraseña.
  \item Ve la lista de aplicaciones.
  \item Puede activar/desactivar prácticas.
  \item Puede programar inicio y fin.
  \item Puede ver informes.
  \item No puede crear usuarios ni cargar HTML.
\end{itemize}
\end{frame}

\begin{frame}{Uso del ESTUDIANTE}
\begin{itemize}
  \item Ingresa con su usuario y contraseña.
  \item Ve solo las prácticas disponibles y autorizadas.
  \item Abre la práctica dentro del sistema.
  \item Puede generar un informe de uso.
  \item Ve solo sus propios registros.
\end{itemize}
\alertbox{El estudiante no tiene funciones de activación, carga o administración.}
\end{frame}

\begin{frame}{Medidas de protección incluidas}
\begin{itemize}
  \item Usuario y contraseña.
  \item Roles verificados en servidor.
  \item Activación/desactivación central.
  \item Horarios de disponibilidad.
  \item HTML almacenado en Drive, no como archivo público directo.
  \item Bloqueo de ejecución local de copias servidas.
  \item Bloqueo disuasorio de clic derecho y atajos comunes.
  \item Aviso de propiedad intelectual e instrucción para IA.
  \item Registro de auditoría e informes.
\end{itemize}
\end{frame}

\begin{frame}{Aviso sobre IA y propiedad intelectual}
\small
El sistema inyecta en cada práctica un aviso que indica que el contenido es privado, restringido y no autorizado para reproducción, extracción o reconstrucción mediante herramientas automatizadas o IA.

\vspace{.3cm}
\begin{block}{Alcance real}
Este aviso sirve como advertencia, evidencia de condiciones de uso y disuasión. No sustituye los controles técnicos.
\end{block}
\begin{block}{Protección más importante}
El HTML se entrega solo después de login y autorización. Las copias locales normales quedan bloqueadas.
\end{block}
\end{frame}

\begin{frame}{Limitaciones honestas}
\begin{itemize}
  \item Ningún sistema web puede impedir absolutamente una foto o captura de pantalla.
  \item Un usuario autorizado puede ver la pregunta que se le muestra.
  \item El bloqueo de clic derecho es disuasorio, no infalible.
  \item Para máxima seguridad, las respuestas correctas deberían evaluarse fuera del HTML original.
\end{itemize}
\alertbox{El objetivo es impedir la distribución libre, controlar accesos y hacer inútiles las copias locales normales.}
\end{frame}

\begin{frame}{Informes y auditoría}
\begin{itemize}
  \item Cada ingreso se registra.
  \item Cada apertura de práctica se registra.
  \item Cada informe generado se registra.
  \item Administradores y docentes ven registros globales.
  \item Estudiantes ven únicamente sus propios registros.
\end{itemize}
\begin{block}{Exportación}
Desde el panel se puede descargar un archivo CSV de informes.
\end{block}
\end{frame}

\begin{frame}{Prueba recomendada antes de usar con estudiantes}
\begin{enumerate}
  \item Crear un usuario DOCENTE de prueba.
  \item Crear un usuario ESTUDIANTE de prueba.
  \item Cargar una aplicación HTML.
  \item Verificar que nace desactivada.
  \item Activarla como DOCENTE.
  \item Entrar como ESTUDIANTE.
  \item Abrir la práctica.
  \item Desactivarla y comprobar que ya no abre.
\end{enumerate}
\end{frame}

\begin{frame}{Problemas frecuentes}
\begin{tabular}{p{4.3cm}p{8.3cm}}
\toprule
\textbf{Problema} & \textbf{Solución}\\
\midrule
No aparece el primer administrador & Verifique que la hoja no tenga usuarios en la pestaña USUARIOS.\\
No carga un HTML & Confirme que el archivo tenga extensión .html o .htm.\\
El estudiante no ve la práctica & Revise que esté activada, dentro de horario y asignada al usuario.\\
Docente no puede crear usuarios & Es correcto: solo ADMINISTRADOR gestiona usuarios.\\
Se pide autorización de Google & Acepte con la cuenta propietaria del sistema.\\
\bottomrule
\end{tabular}
\end{frame}

\begin{frame}{Buenas prácticas}
\begin{itemize}
  \item Use contraseñas de al menos 10 caracteres.
  \item No comparta la cuenta ADMINISTRADOR.
  \item Cree una cuenta DOCENTE para cada profesor responsable.
  \item Desactive prácticas que ya no estén en uso.
  \item Use horarios para evaluaciones o refuerzos temporales.
  \item Descargue informes periódicamente si necesita respaldo externo.
\end{itemize}
\end{frame}

\begin{frame}{Resumen operativo}
\centering
\begin{tikzpicture}[node distance=.7cm, every node/.style={font=\small}]
\tikzstyle{box}=[rounded corners, draw=azul, thick, fill=fondo, align=center, minimum width=8cm, minimum height=.65cm]
\node[box] (a) {Administrador instala una vez};
\node[box, below=of a] (b) {Carga HTML desde el panel};
\node[box, below=of b] (c) {Crea docentes y estudiantes};
\node[box, below=of c] (d) {Docente activa/desactiva con clics};
\node[box, below=of d] (e) {Estudiante ingresa con contraseña};
\node[box, below=of e] (f) {Sistema registra acceso e informes};
\draw[-{Latex}, thick] (a)--(b);\draw[-{Latex}, thick] (b)--(c);\draw[-{Latex}, thick] (c)--(d);\draw[-{Latex}, thick] (d)--(e);\draw[-{Latex}, thick] (e)--(f);
\end{tikzpicture}
\end{frame}

\begin{frame}{Cierre}
\centering
\Large Sistema listo para administración por clics

\vspace{.5cm}
\normalsize
Administrador configura una vez.\par
Docente controla disponibilidad.\par
Estudiante accede solo a lo autorizado.\par

\vspace{.6cm}
\alertbox{Guarde la URL de la aplicación web y la cuenta ADMINISTRADOR inicial.}
\end{frame}

\end{document}
